\documentclass[../../../include/open-logic-section]{subfiles}

\begin{document}

\olfileid{sth}{spine}{Valphabasic}
\olsection{Sipat Dhasar Tataran}

Kanggo nuduhake wigatine definisi himpunan-himpunan
$V_\alpha$ tumrap dhasar teori, kita bakal nerangake
sawetara asil dhasar. Kita miwiti saka sawijining
definisi:\footnote{Durung ana istilah baku kanggo
``potent''; iki jeneng kang dienggo dening
\citet{ButtonLT1}.}
\begin{defn}
	Himpunan $A$ iku \emph{poten} yen lan mung yen
	$\forall x((\exists y \in A)x \subseteq y \lif x \in A)$.
\end{defn}
\begin{lem}\ollabel{Valphabasicprops}
Kanggo saben ordinal $\alpha$:
\begin{enumerate}
	\item\ollabel{Valphatrans} Saben $V_\alpha$ transitif.
	\item\ollabel{Valphapotent} Saben $V_\alpha$ poten.
	\item\ollabel{Valphacum} Yen $\gamma \in \alpha$,
	mula $V_\gamma \in V_\alpha$ (lan mula uga
	$V_\gamma \subseteq V_\alpha$ miturut
	\olref{Valphatrans}).
\end{enumerate}
\end{lem}

\begin{proof}
Kita mbuktekake kanthi induksi transfinit bebarengan.
Kanggo langkah induksi, upamane
\olref{Valphatrans}--\olref{Valphacum} lumaku
kanggo saben ordinal $\beta < \alpha$.

Kasus $\alpha = \emptyset$ cetha.

Upamane $\alpha = \ordsucc{\beta}$. Kanggo mbuktekake
\olref{Valphacum}, yen $\gamma \in \alpha$
mula $V_\gamma \subseteq V_\beta$ miturut
pangira-ira induksi utawa refleksivitas, dadi
$V_\gamma \in \Pow{V_\beta} = V_\alpha$.
Kanggo mbuktekake \olref{Valphapotent}, upamane
$A \subseteq B \in V_\alpha$, yaiku
$A \subseteq B \subseteq V_\beta$; mula
$A \subseteq V_\beta$, dadi $A \in V_\alpha$.
Kanggo mbuktekake \olref{Valphatrans}, yen
$x \in A \in V_\alpha$, kita duwe
$A \subseteq V_\beta$, dadi $x \in V_\beta$.
Mula $x \subseteq V_\beta$ amarga
$V_\beta$ transitif miturut pangira-ira,
lan dadi $x \in V_\alpha$.

Upamane $\alpha$ ordinal limit. Kanggo mbuktekake
\olref{Valphacum}, yen $\gamma \in \alpha$
mula $\gamma \in \ordsucc{\gamma} \in \alpha$.
Mula $V_\gamma \in V_{\ordsucc{\gamma}}$
miturut pangira-ira, saengga
$V_\gamma \in \bigcup_{\beta \in \alpha} V_\beta
= V_\alpha$. Kanggo mbuktekake
\olref{Valphatrans} lan \olref{Valphapotent},
elinga yen gabungan himpunan-himpunan transitif
(utawa poten) tetep transitif (utawa poten).
\end{proof}

\begin{lem}\ollabel{Valphanotref}
Kanggo saben ordinal $\alpha$, $V_\alpha \notin V_\alpha$.
\end{lem}

\begin{proof}
Kanthi induksi transfinit. Cetha yen
$V_\emptyset \notin V_\emptyset$.

Yen $V_{\ordsucc{\alpha}} \in V_{\ordsucc{\alpha}} =
\Pow{V_\alpha}$, mula
$V_{\ordsucc{\alpha}} \subseteq V_\alpha$;
lan amarga $V_\alpha \in V_{\ordsucc{\alpha}}$
miturut \olref{Valphabasicprops}, kita oleh
$V_\alpha \in V_\alpha$. Mula, yen
$V_\alpha \notin V_\alpha$, mesthi
$V_{\ordsucc{\alpha}} \notin V_{\ordsucc{\alpha}}$.

Yen $\alpha$ ordinal limit lan
$V_\alpha \in V_\alpha = \bigcup_{\beta \in
\alpha}V_\beta$, mula $V_\alpha \in V_\beta$
kanggo sawijining $\beta \in \alpha$; nanging
$V_\beta \in V_\alpha$ uga, saengga
$V_\beta \in V_\beta$ miturut
\olref{Valphabasicprops} (dienggo kaping pindho).
Mula, yen $V_\beta \notin V_\beta$ kanggo saben
$\beta \in \alpha$, mesthi
$V_\alpha \notin V_\alpha$.
\end{proof}

\begin{cor}
Kanggo ordinal sebarang $\alpha, \beta$:
$\alpha \in \beta$ yen lan mung yen
$V_\alpha \in V_\beta$.
\end{cor}

\begin{proof}
\Olref{Valphabasicprops} mbuktekake salah siji
arah. Kosokbaline, upamane $V_\alpha \in V_\beta$.
Dadi $\alpha \neq \beta$ miturut
\olref{Valphanotref};
lan $\beta \notin \alpha$, awit yen ora,
kita bakal duwe $V_\beta \in V_\alpha$,
banjur $V_\beta \in V_\beta$ miturut
\olref{Valphabasicprops} (dienggo kaping pindho),
bertentangan karo \olref{Valphanotref}.
Dadi $\alpha \in \beta$ miturut trikotomi.
\end{proof}
\noindent
Kabeh iki ngidini kita nganggep saben $V_\alpha$
minangka tataran kaping-$\alpha$ ing hierarki.
Mangkene alesane.

Saben $V_\alpha$ bisa dianggep kabentuk liwat
proses \emph{iteratif}, amarga panganggone ordinal
nglacak urutan iterasi. Kajaba iku, yen siji
tataran kabentuk luwih dhisik tinimbang tataran
liyane, yaiku $V_\beta \in V_\alpha$, utawa
$\beta \in \alpha$, proses pambentukane
\emph{kumulatif}, amarga
$V_\beta \subseteq V_\alpha$. Pungkasane,
kita pancen mbentuk \emph{kabeh} himpunan kang
bisa dikumpulake saka himpunan-himpunan kang wis
ana ing tataran sadurunge, amarga saben tataran
penerus $V_{\ordsucc{\alpha}}$ iku himpunan
kuasa saka tataran sadurunge, $V_\alpha$.

Dadi, kanthi $\ZFminus$, kita wis
\emph{meh} rampung njlentrehake gagasan babagan
hierarki himpunan kang iteratif lan kumulatif.
(Nanging, kita isih kudu menehi dhasar kanggo
Aksioma Panggantian.)

\end{document}
